\chapter{Spike trains as point processes}\label{ch:points}
\begin{readergoal}Model spike timestamps without pretending they are uniformly sampled voltages.\end{readergoal}
\section{Events rather than amplitudes}
A spike train is a locally finite collection of event times $\{t_1,\ldots,t_K\}$. Its counting process is $N(t)=\#\{j:t_j\le t\}$. It is nondecreasing, integer valued, and jumps at event times. In generalized-function notation $dN(t)/dt=\sum_j\delta(t-t_j)$, but the counting process is the more elementary construction.

\section{Poisson models and their scope}
For a homogeneous Poisson process with rate $\lambda>0$, counts in disjoint intervals are independent and $N(b)-N(a)\sim\mathrm{Poisson}(\lambda(b-a))$. The waiting times are exponentially distributed, and the memoryless model has no refractory period. Real neurons often violate some of these assumptions. A time-inhomogeneous Poisson model admits $\lambda(t)$, with integrated intensity $\int_a^b\lambda(s)ds$ over an interval.

A conditional intensity $\lambda(t\mid\mathcal H_t)$ is a limiting event rate given the observed history $\mathcal H_t$. Roughly,
\[\lambda(t\mid\mathcal H_t)=\lim_{\Delta\downarrow0}\frac{P(N(t+\Delta)-N(t)=1\mid\mathcal H_t)}{\Delta},\]
when the limit is defined and multiple jumps are negligible at small scales. History dependence permits refractory behavior, bursting, and coupling covariates.

\section{Binning and the loss of exact timing}
Dividing time into bins of width $\Delta$ creates counts $n_k=N((k+1)\Delta)-N(k\Delta)$. This may simplify statistics but loses information about positions inside bins. Narrower bins increase temporal precision yet yield sparse counts; wider bins stabilize rate estimates but obscure precise coincidence. A precise synchrony claim therefore needs a time-resolution convention and an explicit null model.

\begin{workedexample}[Poisson probability of a silent window]
A 10-spikes/s homogeneous Poisson process has expected count $\lambda\Delta=0.1$ in 10 ms. The chance of observing zero spikes is $e^{-0.1}\approx0.905$. Frequent silent bins are compatible with a nonzero underlying intensity.
\end{workedexample}
% --- expanded edition (October 2026) ---
\section{Constructing the homogeneous Poisson process}
Three descriptions of the homogeneous Poisson process with rate $\lambda$ are equivalent.
\begin{enumerate}[leftmargin=2em]
\item \emph{Counting description}: $N(0)=0$; counts in disjoint intervals are independent; $N(t+s)-N(t)\sim\mathrm{Poisson}(\lambda s)$.
\item \emph{Interval description}: inter-event intervals $\tau_1,\tau_2,\ldots$ are iid exponential with mean $1/\lambda$, and $t_k=\tau_1+\cdots+\tau_k$.
\item \emph{Infinitesimal description}: $P(N(t+h)-N(t)=1)=\lambda h+o(h)$ and $P(N(t+h)-N(t)\ge2)=o(h)$, with independent increments.
\end{enumerate}
\begin{proposition}
Under the counting description, the first event time is exponential with rate $\lambda$.
\end{proposition}
\begin{proof}
$P(t_1>t)=P(N(t)=0)=e^{-\lambda t}$.
\end{proof}
The interval description gives the simplest simulation: draw $\tau_k=-\log(U_k)/\lambda$ for iid uniform $U_k$ and accumulate. The \emph{conditional uniformity} property is also useful: given $N(T)=n$, the $n$ event times are distributed as the order statistics of $n$ iid uniform points on $[0,T]$. This property underlies spike-time randomization nulls; it holds for homogeneous Poisson processes and fails for any process with refractoriness or bursting.

\section{Inhomogeneous processes, thinning, and time rescaling}
For an inhomogeneous Poisson process with intensity $\lambda(t)$, define the cumulative intensity $\Lambda(t)=\int_0^t\lambda(s)\,ds$. Counts in disjoint intervals are independent and $N(b)-N(a)\sim\mathrm{Poisson}(\Lambda(b)-\Lambda(a))$.

\begin{theorem}[Time rescaling for Poisson processes]
If $t_1<t_2<\cdots$ are the events of an inhomogeneous Poisson process with continuous positive intensity, then $\Lambda(t_1),\Lambda(t_2),\ldots$ are the events of a homogeneous Poisson process with rate one. Equivalently, the rescaled intervals $\Lambda(t_k)-\Lambda(t_{k-1})$ are iid exponential with mean one.
\end{theorem}
\begin{proof}
The map $t\mapsto\Lambda(t)$ is continuous and strictly increasing. The number of rescaled events in $(\Lambda(a),\Lambda(b)]$ equals $N(b)-N(a)$, which is Poisson with mean $\Lambda(b)-\Lambda(a)$, the length of the rescaled interval; independence over disjoint intervals is preserved.
\end{proof}
The general theorem replaces $\lambda(t)$ by the conditional intensity, so that any correctly specified history-dependent model also maps spikes to a unit-rate Poisson process. Testing whether the rescaled intervals look exponential (via $u_k=1-e^{-(\Lambda(t_k)-\Lambda(t_{k-1}))}$, which should be uniform) is the basis of KS plots for spike-train models.

\emph{Thinning} simulates an inhomogeneous process: generate a homogeneous process at rate $\lambda_{\max}\ge\sup_t\lambda(t)$ and keep each event at time $t$ independently with probability $\lambda(t)/\lambda_{\max}$. The kept events form a Poisson process of intensity $\lambda(t)$.

\section{Renewal processes and interval statistics}
A \emph{renewal process} has iid inter-event intervals with a general distribution on $(0,\infty)$. It can express refractoriness (a gap near zero) and regularity, but not dependence between successive intervals. Two summary statistics are standard:
\begin{itemize}[leftmargin=2em]
\item the \emph{coefficient of variation} $\mathrm{CV}=\mathrm{sd}(\tau)/\E(\tau)$, equal to $1$ for the Poisson process, below $1$ for regular firing, and above $1$ for bursty firing;
\item the \emph{Fano factor} $F(T)=\Var N(T)/\E N(T)$ of counts in windows of length $T$.
\end{itemize}
For a stationary renewal process, $F(T)\to\mathrm{CV}^2$ as $T\to\infty$. For finite $T$, both are biased by rate fluctuations across a session, so the local measure $\mathrm{CV}_2=\frac{2|\tau_{k+1}-\tau_k|}{\tau_{k+1}+\tau_k}$, averaged over $k$, is often preferred: it compares adjacent intervals and is insensitive to slow rate drift. A gamma-distributed renewal process with shape $\kappa$ has $\mathrm{CV}=1/\sqrt\kappa$.

\section{Second-order statistics: autocorrelograms and the Bartlett spectrum}
For a stationary point process with rate $\lambda$, the \emph{conditional intensity given an event at $0$}, $\lambda\,g(\tau)$, describes how likely another event is at lag $\tau$. Its histogram estimate is the autocorrelogram. For a homogeneous Poisson process, $g\equiv1$: the autocorrelogram is flat. A dip near zero reflects refractoriness; a peak at $\pm125$\,ms reflects 8\,Hz rhythmicity. The Fourier transform of the covariance density gives the point-process spectrum (Bartlett spectrum), which tends to $\lambda$ at high frequencies for any simple process; a spike train is therefore never ``white'' or ``band-limited'' in the way a continuous signal can be.

\begin{workedexample}[Refractoriness changes the counts]
A neuron firing at 20\,spikes/s with an absolute refractory period of $d=3$\,ms can be modeled as a dead-time-modified Poisson process: intervals are $d+$Exponential$(\mu)$ with $d+1/\mu=50$\,ms. Then $\mathrm{CV}=(1/\mu)/(d+1/\mu)=47/50=0.94$, and the long-window Fano factor is $0.94^2\approx0.88$. Refractoriness alone makes counts sub-Poisson, so a measured Fano factor slightly below one is not evidence of a precise temporal code.
\end{workedexample}
\begin{exercises}
\item Derive the probability of at least one spike in a window under a homogeneous Poisson model.
\item Explain why spike sorting uncertainty can affect autocorrelograms.
\item Distinguish an observed spike count from conditional intensity.
\item Give a point-process model that excludes spikes for a fixed period after each event.
\item Show that the sum of two independent Poisson processes with rates $\lambda_1,\lambda_2$ is a Poisson process with rate $\lambda_1+\lambda_2$.
\item Prove that thinning a homogeneous process with constant retention probability $p$ gives a Poisson process with rate $p\lambda$.
\item Compute $\mathrm{CV}$ for a gamma renewal process with shape $\kappa=4$.
\item Explain why the time-rescaling theorem requires the intensity model to use only past information.
\end{exercises}
